% resultados.tex
\section{Resultados}\label{sec:resultados}

\subsection{Implementación Práctica}

Los algoritmos fueron implementados en MATLAB y probados con imágenes sintéticas y reales. 
Las implementaciones están disponibles en el apéndice de este documento.

\subsection{Complejidad Computacional}

\begin{table}[h]
\centering
\caption{Comparación de complejidad de métodos}
\begin{tabular}{lccc}
\hline
\textbf{Método} & \textbf{Operaciones por pixel} & \textbf{Memoria} & \textbf{Tiempo relativo} \\
\hline
Lineal & 2-4 & Baja & 1.0× \\
Bicuadrada & 8-12 & Media & 2.5× \\
Bicúbica & 16-24 & Alta & 4.0× \\
\hline
\end{tabular}
\end{table}

\subsection{Calidad de Interpolación}

La evaluación de calidad mostró que la interpolación bicúbica proporciona los mejores resultados visuales, seguida por la bicuadrada y finalmente la lineal. La interpolación bicúbica preserva mejor los bordes y detalles finos, mientras que la lineal tiende a producir artefactos de aliasing en regiones de alto contraste. La bicuadrada ofrece un compromiso aceptable entre calidad y complejidad computacional.

\subsection{Problemas en Bordes}

La interpolación en los bordes de la imagen presenta desafíos particulares, ya que no existen suficientes vecinos para aplicar los métodos estándar. Se requieren técnicas especiales como extrapolación o espejado para manejar estas regiones correctamente. Todos los métodos evaluados requieren tratamiento especial en los bordes, siendo más crítico en los métodos de mayor orden que utilizan vecindarios más grandes.

\subsection{Vecindarios Utilizados}

\begin{table}[h]
\centering
\caption{Configuraciones de vecindario para interpolación}
\begin{tabular}{lp{3cm}p{3cm}p{3cm}}
\hline
\textbf{Método} & \textbf{Patrón 1D} & \textbf{Tamaño} & \textbf{Puntos utilizados} \\
\hline
Lineal & [P₁ ? P₂] & 3×1 & 2 \\
Bicuadrada & [P₁ P₂ ? P₃ P₄] & 5×1 & 3 \\
Bicúbica & [P₁ P₂ P₃ ? P₄ P₅ P₆] & 7×1 & 4 \\
\hline
\end{tabular}
\end{table}